\section{special functions}

It is certainly very useful to know the methods of integration in order to
explicitly write the antiderivative of a generic function. However, it is also
very useful to know that there exist some elementary functions whose
antiderivative is not an elementary function. A prime example is the
\emph{Gaussian bell curve}%
\mymargin{Gaussian bell curve}%
\index{bell curve!Gaussian}%
\index{Gauss!function of}%
\index{function!Gaussian}%
$f(x) = e^{-x^2}$ whose integral function 
\[
 F(x) = \int_0^x e^{-t^2}\, dt
\]
cannot be expressed through the composition of elementary functions. This fact
should not discourage us too much; we will assign a name to the function $F$ and
study its properties using the theory we have developed. Obviously, it is likely
that someone before us has encountered such functions, already named them, and
studied their properties. These functions are usually called \emph{special
functions} to distinguish them from \emph{elementary functions}.
\index{function!special}%
\index{special!function}%
\index{functions!elementary}%
\index{elementary!function}%
In the specific example, the \emph{error function}%
\mymargin{error function}%
\index{function!error}
has been defined as
\[
  \erf x = \frac{2}{\sqrt \pi} \int_0^x e^{-t^2}\, dt.
\]

Other examples of functions whose antiderivative cannot be expressed using
elementary functions are the following:
\[
  \frac{1}{\ln x}, \qquad 
  \frac{e^x}{x}, \qquad
  \frac{\sin x}{x}, \qquad \sin\frac{\pi x^2}{2}
\]
for which the respective antiderivatives are defined: \emph{logarithmic
integral}, \emph{exponential integral}, \emph{sinusoidal integral},
\emph{Fresnel integral}:
\[
\li x, \qquad 
\ei x, \qquad 
\Si x, \qquad
S x.
\]

The theorem through which it is demonstrated that these functions do not admit
an antiderivative expressible as a composition of elementary functions (rational
functions, exponential, trigonometric, and their inverses) is called Liouville's
Theorem.%
\newsavebox{\qrDeLellis}\sbox{\qrDeLellis}{%
\myqrdoclink{http://cvgmt.sns.it/paper/3456/}{}{Liouville's Theorem}}%
\mynote{%
Liouville's Theorem is an algebraic type theorem and thus goes far beyond the
scope of this course. For those interested in seeing its proof, a pleasant read
can be the article by Camillo De Lellis: \emph{Liouville's Theorem or why
``there does not exist'' the antiderivative of $\exp(x^2)$}.\\
\usebox{\qrDeLellis}
}